\subsection{A polynomial finite-prefix certificate under a denominator bound}

\begin{proposition}[P-F13: universal denominator and finite reconstruction]
\label{prop:finiteprefix}
Let $R,A\ge1$ be given integers. Suppose a nonzero global
quasi-polynomial $f(n)$, $n\ge0$, has a positive rational form
\[
 F(T)=\sum_{n\ge0}f(n)T^n
 =\frac{N(T)}{Q(T)},\qquad
 Q(T)=\prod_{j=1}^{r}(1-T^{a_j}).
\]
\[
 1\le r\le R,\quad 1\le a_j\le A,
 \quad N\in\mathbb Z_{\ge0}[T]\setminus\{0\}.
\]
Define
\[
 D(T)=\prod_{a=1}^{A}(1-T^a)^R,
 \qquad M=\deg D=\frac{RA(A+1)}2.
\]
Then $F=P/D$ for a unique polynomial $P\in\mathbb Z[T]$ with
$\deg P<M$, and
\[
 P(T)=D(T)\left(\sum_{n=0}^{M-1}f(n)T^n\right)\bmod T^M.
 \tag{P-F13a}
\]
Consequently those first $M$ coefficients determine the entire
generating function. Under the stated denominator promise, a proposed
positive form with at most $R$ denominator factors of weights at most
$A$ and with numerator degree smaller than its denominator degree can
be certified by these $M$ coefficients.

Let $b$ bound the binary lengths of the supplied integers
$f(0),\ldots,f(M-1)$. Reconstruction and verification of any proposed
denominator take bit complexity polynomial in the numeric parameters
$R,A$ and $b$. A positive form can be found by a finite search through
at most $\binom{A+R}{R}-1$ denominators, with polynomial work per
denominator, once those coefficient values are supplied.
Equivalently, positive-form selection for the reconstructed rational
function has certificates verifiable in polynomial time in this data.
No polynomial-time algorithm for evaluating the original multiplicities
or selecting a successful denominator is assumed or proved.
\end{proposition}

\begin{proof}
Let $c_a$ be the number of occurrences of $a$ among the $a_j$.
Then $0\le c_a\le R$, and
\[
 D/Q=\prod_{a=1}^{A}(1-T^a)^{R-c_a}\in\mathbb Z[T].
\]
By the global quasi-polynomial properness argument in P-F9,
$\deg N<\deg Q$. Hence
\[
 P=N(D/Q)\in\mathbb Z[T],\qquad
 \deg P=\deg N+M-\deg Q<M.
\]
Since $D(0)=1$, multiplying the first $M$ coefficients of $F$ by $D$
and discarding terms of degree at least $M$ recovers every coefficient
of $P$, proving (P-F13a). Uniqueness follows immediately. More
explicitly, two strictly proper rational functions whose denominators
divide $D$ and whose first $M$ coefficients agree have a difference
$(P_1-P_2)/D$ with $\deg(P_1-P_2)<M$ and zero of order at least $M$
at zero. Multiplication by $D(0)\ne0$ forces $P_1-P_2=0$.

For an effective verifier, fix a candidate list
$1\le a_1\le\cdots\le a_r\le A$, $1\le r\le R$, and put
$Q'=\prod_j(1-T^{a_j})$ and $d'=\deg Q'\le RA\le M$.
Its only possible strictly proper numerator is
\[
 N'(T)=Q'(T)\left(\sum_{n=0}^{M-1}f(n)T^n\right)\bmod T^{d'}.
 \tag{P-F13b}
\]
Accept this list exactly if all coefficients of $N'$ are nonnegative,
$N'\ne0$, and the polynomial identity
\[
 D(T)N'(T)=Q'(T)P(T)
 \tag{P-F13c}
\]
holds. If it is accepted, (P-F13c) proves $N'/Q'=P/D=F$.
Conversely any positive form with the stated bounds has a strictly
proper numerator, is recovered by (P-F13b), and passes (P-F13c).
Thus the test is both necessary and sufficient under the promise.
The same argument proves the first-$M$-coefficients version of the
certificate.

Here is an explicit bit bound. Write $\|B\|_1$ for the sum of the
absolute values of the coefficients of a polynomial $B$. One has
$\|D\|_1\le2^{AR}$ and $\|Q'\|_1\le2^R$. If
$|f(n)|<2^b$ for $n<M$, formula (P-F13a) bounds the absolute value
of every coefficient of $P$ by $2^{b+AR}$, and (P-F13b) bounds that
of every coefficient of $N'$ by $2^{b+R}$. All degrees in the
construction and in (P-F13c) are at most $M+RA$. Ordinary integer
polynomial multiplication, truncation and equality tests therefore
have bit complexity polynomial in $R,A,b$. This counts $R,A$ by
their numeric values; it does not claim polynomial complexity in
$\log R,\log A$.

There are $\binom{A+r-1}{r}$ weakly increasing lists of length $r$,
so the sum for $1\le r\le R$ is $\binom{A+R}{R}-1$.
Enumerating these and using the verifier must find an accepted list
under the assumed existence of a positive form. The list, together
with the deterministically recovered numerator if desired, is a
polynomial-size witness for successful selection from the supplied
finite data. The number of lists is not asserted to be polynomial.
\end{proof}

\paragraph{Application and scope.}
Under a specified bound $A\le A_0$ in PH3, P-F12 supplies a
representation-theoretic bound $r\le R$. P-F13 reduces construction of
a positive form to the first $RA_0(A_0+1)/2$ exact stretching
multiplicities and finite denominator selection. For $G=GL_m$, the
Weyl estimate in P-F12 bounds the binary lengths of all these values
by a polynomial in $m,R,A_0,\log(\lambda_1+1)$. The number of
required coefficients and the work of verifying a candidate are
therefore polynomial in these expanded parameters. This avoids
verification across a possibly exponential least-common-multiple
period, and it does not provide an efficient coefficient oracle.

The denominator promise is essential. A finite prefix, by itself,
does not prove PH3 for an otherwise unbounded original series.
For example, the two nonnegative global quasi-polynomials
$1$ and $1+\mathbf1_{\{n\equiv M\pmod{M+1}\}}$ agree for
$n=0,\ldots,M-1$ and differ at $n=M$. Thus (P-F13c) certifies the
original multiplicity series only when its membership in the stated
denominator class is known, as it is conditionally under PH3.
Historical priority for this elementary rational reconstruction is
not asserted.
